\begin{definition}
    Sea $x=(x_n)_{n\geq1}$ una sucesión de números reales. Decimos que $(x_n)$ es
    equidistribuida módulo $1$ si para cada
    intervalo $(a,b) \subset \mathbb{T}$ con $0 < a < b < 1$ se cumple que
    \begin{align}
        \label{def:equidistribucion}
        \lim_{N\to\infty} \frac{A_N((a,b), x)}{N} = b-a
    \end{align}
    donde 
    \[
        A_N((a,b), x) = \card \left\{1\le n\le N : \partfrac{x_n} \in (a,b)\right\}
    \]
\end{definition}

\begin{remark}
    \label{obs:equidistribucion}
    Si consideramos la función característica $\mathbf{1}_{(a,b)}$ en $\mathbb{T}$, entonces la
    condición de equidistribución se puede reescribir como
    \[
        \lim_{N\to\infty} \frac{1}{N} \sum_{n=1}^N \mathbf{1}_{(a,b)}(x_n) =
        \int_0^1 \mathbf{1}_{(a,b)}(x)\ dx = 
        b-a
    \]
\end{remark}